Visual programming languages are widely used to introduce young learners to
programming, but it is less clear how well they support the next step:
teaching software design principles. This study compares two visual programming
paradigms in that context, block-based with Scratch and flow-based with Godot
Orchestrator, and investigates how each conveys modularity, instantiation,
state management and event handling to young learners. We conducted a two-day
workshop with three middle-school students, in which two groups built the
same game using either of the two languages. Data was collected through pre-
and post-surveys, a multiple-choice test, observations during the workshop,
and evaluation of the student-created artifacts. Our results indicate that
Scratch correlated with a better understanding of the selected software
design principles than Godot Orchestrator, especially for modularity, while
instantiation turned out to be the hardest concept to grasp in both
languages. At the same time, observations suggest that the Scratch student
overestimated his abilities, and that Scratch's high level of abstraction
limits teaching more advanced concepts such as finite state machines. The
Orchestrator students progressed slower, but showed substantial improvement
in their use of variables and functions over the course of the workshop.
Beyond the language comparison, we find that familiarity with the chosen
language might matter more than its specific paradigm when teaching software
design, as students new to a language spent considerable time learning the
tool rather than the concepts. Due to the low number of participants, this
study serves as a pilot rather than a conclusive comparison, but it points
towards interesting directions for future research on teaching software
design with visual programming languages.
